% Source: https://artofproblemsolving.com/community/c3h2594653_trig_questions_that_i_cant_understand_the_solutions_to_
% Stable identity: aops-topic-2594653-q03
\textbf{(a)} Over the real numbers, $\arcsin x$ is defined only for $-1\le x\le1$: it is the angle in $\left[-\dfrac\pi2,\dfrac\pi2\right]$ whose sine is $x$. So $\sin(\arcsin x)$ is defined exactly for $-1\le x\le1$, and then it equals $x$; for $|x|>1$ it is undefined. In the same way $\arccos x$ is the angle in $[0,\pi]$ whose cosine is $x$, so $\cos(\arccos x)=x$ for $-1\le x\le1$ and is undefined otherwise. In short, both are $\boxed{x\text{ for }|x|\le1,\text{ undefined for }|x|>1}$.

\textbf{(b)} $\arcsin(\sin x)$ is defined for every real $x$, but its value always lies in $\left[-\dfrac\pi2,\dfrac\pi2\right]$. If $x$ is in this interval, it is the angle there with sine $\sin x$, so $\arcsin(\sin x)=x$. Otherwise the value differs from $x$. So $\arcsin(\sin x)=x$ exactly for $\boxed{-\dfrac\pi2\le x\le\dfrac\pi2}$.

\begin{asy}
usepackage("amsmath");
size(7cm);
defaultpen(fontsize(10pt));
import graph;
real f(real x){return asin(sin(x));}
draw((-7,0)--(7,0),Arrow(4));draw((0,-2.2)--(0,2.4),Arrow(4));
draw(graph(f,-2*pi,2*pi,400),linewidth(0.9));
draw((-1.9,-1.9)--(1.9,1.9),dashed);
draw((pi/2,0)--(pi/2,pi/2),dotted);draw((-pi/2,0)--(-pi/2,-pi/2),dotted);
label("$\dfrac\pi2$",(pi/2,0),S,fontsize(8pt));label("$-\dfrac\pi2$",(-pi/2,0),N,fontsize(8pt));
label("$y=\arcsin(\sin x)$",(2*pi,0),N,fontsize(8pt));label("$y=x$",(1.9,1.9),NE,fontsize(8pt));
\end{asy}

The graph of $y=\arcsin(\sin x)$ is a zigzag that agrees with $y=x$ only between $-\dfrac\pi2$ and $\dfrac\pi2$.
